% Source slides 77--80
\section[Partitioned matrices]{Partitioned Matrices}
\sectioncard{2.4}{Partitioned Matrices\\[-.25ex]{\Large 分块矩阵}}

\begin{frame}{Row Partitions}
\centering
\[
A=\left[
\begin{array}{cccc}
a&1&0&0\\ \hdashline
0&a&0&0\\
1&0&b&1\\ \hdashline
0&1&1&b
\end{array}\right]
=\begin{bmatrix}B_1\\B_2\\B_3\end{bmatrix}.
\]
\medskip
The horizontal cuts may contain any number of rows. For example,
\[
A=\left[
\begin{array}{cccc}
a&1&0&0\\ \hdashline
0&a&0&0\\
0&1&1&b\\ \hdashline
0&1&1&b
\end{array}\right]
=\begin{bmatrix}B_1\\B_2\\B_3\end{bmatrix}.
\]
\end{frame}

\begin{frame}{Block Partitions}
\centering
\[
A=\left[
\begin{array}{cc:cc}
a&1&0&0\\
0&a&0&0\\ \hdashline
1&0&b&1\\
0&1&1&b
\end{array}\right]
=\begin{bmatrix}C_1&C_2\\C_3&C_4\end{bmatrix}.
\]
\medskip
Each $C_i$ is itself a matrix. Here every block is $2\times2$.
\end{frame}

\begin{frame}{Common Partition Patterns}
Let
\[
A=\begin{bmatrix}
a&1&0&0\\
0&a&0&0\\
1&0&b&1\\
0&1&1&b
\end{bmatrix}.
\]
Then two useful descriptions are
\[
A=\begin{bmatrix}D&O\\I&B\end{bmatrix},\qquad
D=\begin{bmatrix}a&1\\0&a\end{bmatrix},\quad
B=\begin{bmatrix}b&1\\1&b\end{bmatrix},
\]
and the column partition
\[
A=[\,A_1\ A_2\ A_3\ A_4\,],\qquad
A_2=\begin{bmatrix}1\\a\\0\\1\end{bmatrix}.
\]
\end{frame}
